\subsection{Power Functions} 
For each possible exponent $p$, we have a family of \term{power functions}.
\begin{definition}[Power Function]
A \term{power function} is one that can be written in the form
$$f(x)=a\left(bx-h\right)^p+k$$
where $a$, $b$, $h$, and $k$ are any real numbers. This is a family of functions with base function
$$f(x)=x^p$$
Each exponent $p$ gives a different family.
\end{definition}
\noi In practice, we will usually exclude exponents \makeblanksmall and \makeblanksmall:
\begin{itemize}
\item If $p=\makeblanksmall$, the base function is $f(x)=\makeblank[1in]$, which is a \makeblank[1in] function.
\item If $p=\makeblanksmall$, the base function is $f(x)=\makeblank[1in]$, which is a \makeblank[1in] function.
\end{itemize}
To begin, we will limit our discussion to \makeblank exponents.